\section{Discussion}

The synthesis of the six themes supports four transversal readings that cut across the research questions and organize what the corpus, taken as a whole, tells us about the state of the field.

\noindent\textbf{The unit of work has shifted from prompt to process.} The central thesis of this study holds empirically. In 34 of the 37 studies the contribution is described as a framework, and the fields for persistent artifacts, roles, and validation are present in the vast majority of records. The language model is rarely the declared differentiator; what sets the proposals apart is the process structure that organizes specs, roles, execution, and validation around the agent's work. The three studies that are not frameworks reinforce the contrast: an architecture of prompt patterns, an AI-chain platform, and a method of future scenarios. These are exactly the cases where the unit is still the isolated technique and not the process \cite{JulesWhite2023,YuCheng2023,JSauvola2024}.

\noindent\textbf{The six-dimension taxonomy is partly validated and partly challenged.} Five dimensions (specification, context, roles, execution, and validation) have strong empirical anchoring in at least one theme. The sixth, portability, is weak and almost always diagnostic: it shows up as a risk of lock-in and platform dependence, not as a deliberate design property. Two relevant axes also emerge outside the six dimensions, the human-organizational one in T5 and the security one in T6. The analytical recommendation that follows from the hybrid mode is concrete: consider extending the taxonomy with a dimension of human-AI collaboration and governance, and treat security as a transversal evaluation criterion or a subdimension of validation. That is what confronting the themes with the taxonomy adds, compared to assuming the dimensions a priori.

\noindent\textbf{Three recurring axes of tension.} The corpus shows partial consensus, not convergence, along three axes. The first is autonomy versus human control: CodePori, MetaGPT, and ChatDev lean toward an end-to-end autonomous flow, while frameworks centered on adoption and trust build human decision in as a design element, and the distinction between vibe coding and agentic coding theorizes the contrast \cite{ZeeshanRasheed2024,SiruiHong2023,ChenQian2023,SaiZhang2025,DanielRusso2024,Baron2025,RanjanSapkota2025}. The second is heavy governance versus lightweight agility: verification ceremonies and formal contracts stand against lean flows for startups and prompt patterns, and coordination cost is the recurring price of governance \cite{Merchant2025,Paduraru2026,Khan2026,JulesWhite2023}. The third is the direction of the specification flow: specifying first, in spec-driven development, MetaGPT, and AgileGen, versus recovering the specification from legacy, in Reversa \cite{DeepakBabuPiskala2026,SiruiHong2023,SaiZhang2025,Macedo2026}. The two extremes anchor in the same specification dimension but in opposite directions, which is material for any future comparative matrix.

\noindent\textbf{Traceability and context are the most cited transversal requirements.} The risk tied to context appears in nearly every study, and traceability is a recurring demand in the specification and governance themes. Together, they indicate that the field's practical bottlenecks are not in code generation itself, but in keeping the agent anchored in the real project and auditable over time. This finding realigns the usual reading of the problem: the open challenge is less about producing plausible code and more about sustaining fidelity to the system and to the recorded intent.

\noindent\textbf{The evidence is uneven.} Although 26 of the 37 studies declare empirical evaluation and 22 have a public repository, ten are arXiv preprints not yet peer-reviewed, and most evaluations occur on benchmarks or one-off case studies, not on longitudinal industrial projects. The strength of the evidence varies widely by source category, exactly as the protocol anticipated by classifying evidence by category instead of applying quality-based exclusion. This heterogeneity calls for caution: the findings of this study describe the structure and tensions of the field with confidence, but claims about relative performance between frameworks remain underdetermined by the available evidence.
